\RequirePackage{plautopatch}
\documentclass[14pt,aspectratio=169,xcolor=dvipsnames,table,onlytextwidth,dvipdfmx]{beamer}
\usepackage[utf8]{inputenc}
\usepackage{bxdpx-beamer} % dvipdfmxなので必要

\usetheme{Boadilla}
%Beamerフォント設定
\usefonttheme{professionalfonts}
\usepackage[T1]{fontenc}
\usepackage[deluxe,uplatex]{otf} % 日本語多ウェイト化
\usepackage{mlmodern}  % 太いComputer Modern
% MLmodernのバグを修正: cf. https://tex.stackexchange.com/questions/646333/size-of-integral-symbol-in-section-header-with-mlmodern
\DeclareFontFamily{OMX}{mlmex}{}
\DeclareFontShape{OMX}{mlmex}{m}{n}{%
   <->mlmex10%
   }{}
\usepackage{newtxtext} % 数式以外をTXフォントで上書き
% \usepackage[noalphabet,unicode]{pxchfon}
% \setminchofont{KozMinPro-Regular.otf}% 游明朝Regular
% \setboldminchofont{KozMinPro-Bold.otf}% 游明朝Demibold
% \setgothicfont[0]{BIZ-UDGothicR.ttc}% BIZ UDゴシックR
% \setboldgothicfont[0]{BIZ-UDGothicB.ttc}% BIZ UDゴシックB
\renewcommand{\familydefault}{\sfdefault}  % 英文をサンセリフ体に
\renewcommand{\kanjifamilydefault}{\gtdefault}  % 日本語をゴシック体に
\usefonttheme{structurebold} % タイトル部を太字
\setbeamerfont{alerted text}{series=\bfseries} % Alertを太字
\setbeamerfont{section in toc}{series=\mdseries} % 目次は太字にしない
\setbeamerfont{frametitle}{size=\Large} % フレームタイトル文字サイズ
\setbeamerfont{title}{size=\LARGE} % タイトル文字サイズ
\setbeamerfont{date}{size=\small}  % 日付文字サイズ
\usepackage{bxcoloremoji}

% Babel (日本語の場合のみ・英語の場合は不要)
\uselanguage{japanese}
\languagepath{japanese}
\deftranslation[to=japanese]{Theorem}{定理}
\deftranslation[to=japanese]{Lemma}{補題}
\deftranslation[to=japanese]{Example}{例}
\deftranslation[to=japanese]{Examples}{例}
\deftranslation[to=japanese]{Definition}{定義}
\deftranslation[to=japanese]{Definitions}{定義}
\deftranslation[to=japanese]{Problem}{問題}
\deftranslation[to=japanese]{Solution}{解}
\deftranslation[to=japanese]{Fact}{事実}
\deftranslation[to=japanese]{Proof}{証明}
\deftranslation[to=japanese]{Corollary}{系}
\setbeamertemplate{theorems}[normal font] %日本語の場合，定理内を斜体にしない
% overlay-aware proof environment
\renewenvironment<>{proof}{%
\par
\begin{actionenv}#1%
\structure{(証明)}}
{\qed\end{actionenv}}

%Beamer色設定
\definecolor{UniBlue}{RGB}{0,150,200}
\definecolor{UniGreen}{RGB}{3,175,122}
\definecolor{UniPink}{RGB}{255,128,255}
\definecolor{AlertOrange}{RGB}{255,76,0}
\definecolor{AlmostBlack}{RGB}{38,38,38}
\setbeamercolor{normal text}{fg=AlmostBlack}  % 本文カラー
\setbeamercolor{structure}{fg=UniBlue} % 見出しカラー
\setbeamercolor{block title}{fg=UniBlue!50!black} % ブロック部分タイトルカラー
\setbeamercolor{alerted text}{fg=AlertOrange} % \alert 文字カラー
\setbeamercolor{bibliography entry author}{fg=AlmostBlack} % Bib 文字カラー
\setbeamercolor{bibliography entry note}{fg=AlmostBlack} % Bib 文字カラー
\mode<beamer>{
    \definecolor{BackGroundGray}{RGB}{254,254,254}
    \setbeamercolor{background canvas}{bg=BackGroundGray} % スライドモードのみ背景をわずかにグレーにする
}

%フラットデザイン化
\setbeamertemplate{blocks}[rounded] % Blockの影を消す
\useinnertheme{circles} % 箇条書きをシンプルに
\setbeamertemplate{navigation symbols}{} % ナビゲーションシンボルを消す
\setbeamertemplate{footline}[frame number] % フッターはスライド番号のみ

%タイトルページ
\setbeamertemplate{title page}{%
    \vspace{2.5em}
    {\usebeamerfont{title} \usebeamercolor[fg]{title} \inserttitle \par}
    {\usebeamerfont{subtitle}\usebeamercolor[fg]{subtitle}\insertsubtitle \par}

    \vspace{1.5em}
    \begin{columns}[]
        \begin{column}{.35\linewidth}
        \usebeamerfont{titlegraphic}\inserttitlegraphic\par
        \end{column}
        \begin{column}{.6\linewidth}\setlength\topsep{0pt}
            \begin{flushright}
        \usebeamerfont{author}\insertauthor\par
        \usebeamerfont{institute}\insertinstitute \par
        \vspace{3em}
        \usebeamerfont{date}\insertdate\par
            \end{flushright}
        \end{column}
    \end{columns}
}

%biblatex
% \usepackage[style=ext-authoryear,citestyle=authoryear,natbib=true,giveninits=true,maxbibnames=99,maxcitenames=10,url=false,isbn=false,doi=false,articlein=false]{biblatex}
% \DeclareNameAlias{author}{first-last}
% \addbibresource{shrunk.bib}
% \newcommand{\mkbibbracketscol}[1]{\textcolor{UniBlue}{\scriptsize \mkbibbrackets{#1}}}
% \DeclareCiteCommand{\cite}[\mkbibbracketscol]
%   {\usebibmacro{prenote}}
%   {\usebibmacro{citeindex}%
%    \usebibmacro{cite}}
%   {\multicitedelim}
%   {\usebibmacro{postnote}}
% % \renewcommand{\cite}[1]{\textcolor{UniBlue}{\scriptsize\citep{#1}}}
% \renewcommand*{\finalnamedelim}{\addcomma\addspace} % remove "and" before the last author
% \DeclareDelimFormat{nameyeardelim}{\addspace}
% \setbeamertemplate{bibliography item}{}

% Algorithm系
\usepackage{algorithm}
\usepackage[noend]{algorithmic}
\algsetup{linenosize=\color{fg!50}\footnotesize}
\renewcommand\algorithmicdo{:}
\renewcommand\algorithmicthen{:}
\renewcommand\algorithmicrequire{\textbf{Input:}}
\renewcommand\algorithmicensure{\textbf{Output:}}
\renewcommand\algorithmiccomment[1]{\hfill\textcolor{fg!50}{\footnotesize\textit{// #1}}}

% TikZ
\usepackage{tikz}
\usetikzlibrary{positioning,shapes,arrows,quotes,shapes.callouts,calc,graphs,graphs.standard,fit,backgrounds,perspective,matrix}
\tikzset{point/.style={circle, fill=AlertOrange, inner sep=1pt, text width=1pt}}
\tikzset{bpoint/.style={circle, fill=black, inner sep=1pt, text width=1pt}}
\tikzset{alt/.code args={<#1>#2#3}{\alt<#1>{\pgfkeysalso{#2}}{\pgfkeysalso{#3}}}}
\tikzset{>=latex}
\tikzset{mynodes/.style={circle,white,fill=UniBlue!90,text width=.8em,inner sep=0pt,text centered,font=\footnotesize}}
\tikzset{myedges/.style={thick}}
\tikzset{mypath/.style={ultra thick, draw=#1}, mypath/.default=AlertOrange}
\tikzset{mysubset/.style={draw,#1,dotted,thick,fill=#1!10}, mysubset/.default=UniBlue}
\tikzset{mypath2/.style={ultra thick, dashed, draw=UniGreen}}
\tikzset{graphs/mygraph/.style={nodes={mynodes}, edges={myedges}, empty nodes}}
\tikzset{nomargin/.style={inner sep=0pt, outer sep=0pt}}
%% graph macros
\tikzgraphsset{declare={mybipartite}{%
subgraph I_nm [V={1,2,3,4,5,6}, W={a,b,c,d,e,f}];
1 -- {a,b};
2 -- {b,c,d,e};
3 -- {d,f};
4 -- {e,f};
{5,6} -- f;
}}
\tikzgraphsset{declare={mybipartitedi}{%
subgraph I_nm [V={1,2,3,4,5,6}, W={a,b,c,d,e,f}];
1 -> {a,b};
2 -> {b,c,d,e};
3 -> {d,f};
4 -> {e,f};
{5,6} -> f;
}}
\tikzgraphsset{declare={mybipartite2}{%
subgraph I_nm [V={1,2,3}, W={a,b,c}];
1 -- {a,b}; 2 -- {a,b,c}; 3 -- {a,c};
}} 
% Petersen graph
\tikzgraphsset{declare={Petersen}{%
    subgraph C_n [n=5,name=A, radius=1.5cm]; 
    subgraph I_n [n=5,name=B, radius=.8cm];
    \foreach \i [evaluate={\j=int(mod(\i+2,5)+1);}] in {1,...,5}{
        A \i -- B \i;
        B \i -- B \j;
}}}
% Example graph for Menger's theorem
\tikzgraphsset{declare={menger}{%
    s[at={(-1,0)}, as={$s$}];
    (s) -> 1[at={(0,1)}] -> 2[at={(1,0)}];
    (s) -> 3[at={(0,-1)}] -> 2;
    4[at={(2,1)}] <- 5[at={(2,-1)}];
    {5,6[at={(3,1)}], 7[at={(3,-1)}]} -> t[at={(4,0)}, as={$t$}];
    (s) -> (2);
    (2) -> {(4), (5)};
    {(4),(5)} -> {(6),(7)};
    (6) -> (7);
    1 -> 4;
}} 
\tikzgraphsset{declare={undirected menger}{%
    s[at={(-1,0)}, as={$s$}];
    (s) -- 1[at={(0,1)}] -- 2[at={(1,0)}];
    (s) -- 3[at={(0,-1)}] -- 2;
    4[at={(2,1)}] -- 5[at={(2,-1)}];
    {5,6[at={(3,1)}], 7[at={(3,-1)}]} -- t[at={(4,0)}, as={$t$}];
    (s) -- (2);
    (2) -- {(4), (5)};
    {(4),(5)} -- {(6),(7)};
    (6) -- (7);
    1 -- 4;
}} 

% Appendix
\usepackage{appendixnumberbeamer}

% macros
\newcommand{\R}{\mathbb{R}}
\newcommand{\Z}{\mathbb{Z}}
\newcommand{\be}{\mathbf{e}}
\newcommand{\zeros}{\mathbf{0}}
\newcommand{\ones}{\mathbf{1}}
\newcommand{\bfzero}{\zeros}
\newcommand{\bfone}{\ones}

% \newcommand{\bp}{\mathbf{p}}
\newcommand{\eps}{\varepsilon}
\newcommand{\defiff}{\overset{\text{def}}{\iff}}
\DeclareMathOperator{\poly}{poly}
\DeclareMathOperator{\conv}{conv}
\DeclareMathOperator{\In}{In}
\DeclareMathOperator{\Out}{Out}
\DeclareMathOperator{\val}{val}
\newcommand{\lo}{{\mathrm{lo}}}
\newcommand{\up}{{\mathrm{up}}}
\newcommand{\caI}{\mathcal{I}}
\newcommand{\caB}{\mathcal{B}}
\newcommand{\caC}{\mathcal{C}}
\newcommand{\bM}{\mathbf{M}}
\newcommand{\IO}{\mathrm{IO}}

\usepackage{mathtools}
\DeclarePairedDelimiter{\abs}{\lvert}{\rvert}
\DeclarePairedDelimiter{\norm}{\lVert}{\rVert}
\DeclarePairedDelimiter{\inprod}{\langle}{\rangle}
\usepackage[no-test-for-array]{nicematrix}
\usepackage{diagbox}

\newcommand{\highlightcap}[3][red]{\tikz[baseline=(x.base)]{\node[rectangle,rounded corners,fill=#1!10](x){#2} node[below=0pt of x, color=#1]{#3};}}
\newcommand{\highlight}[2][red]{\tikz[baseline=(x.base)]{\node[rectangle,rounded corners,fill=#1!10](x){#2};}}
\newcommand{\mycite}[2][\footnotesize]{\textcolor{UniBlue}{#1 [#2]}}

%section
% \AtBeginSection[]{
%     \frame[c]{
%         \centering
%         \begin{tikzpicture}
%             \node[white, circle, fill=UniBlue, text centered, text width=.3\linewidth](bullet) {
%                 {\Huge \bf \insertsectionnumber} \\[4em]
%     };
%     \node[text width=.65\linewidth, right=1em of bullet]{%
%         \usebeamerfont{title}\usebeamercolor[fg]{title}\insertsection\par %
%     };
%         \end{tikzpicture}
%     } %目次スライド
% }

%grid
% \usepackage[colorgrid,gridunit=pt,texcoord]{eso-pic}
\newcommand{\postit}[1]{\tikz[overlay, remember picture]{\node[draw,anchor=north east, align=center, rectangle, fill=yellow!10, xshift=-1pt, yshift=-1pt] at (current page.north east) {#1};}}

\usepackage{tcolorbox}
\newtcolorbox{primal}{colback=red!10!white, colframe=red, title=主問題 (P), left=2pt,right=2pt,top=0pt,bottom=0pt}
\newtcolorbox{dual}{colback=blue!10!white, colframe=blue, title=双対問題 (D), left=2pt,right=2pt,top=0pt,bottom=0pt}

% visible on
\usetikzlibrary{overlay-beamer-styles}
\usepackage[normalem]{ulem}
\usepackage{pxrubrica}
\usepackage{booktabs}

\newcommand{\Aphiup}{\textcolor{UniBlue}{A_\varphi^\up}}
\newcommand{\Aphilow}{\textcolor{AlertOrange}{A_\varphi^\lo}}
\newcommand{\bv}{\boldsymbol{v}}
\newcommand{\bx}{\boldsymbol{x}}
\newcommand{\by}{\boldsymbol{y}}
\newcommand{\bz}{\boldsymbol{z}}
\newcommand{\bw}{\boldsymbol{w}}
\newcommand{\bc}{{\boldsymbol c}}
\DeclareMathOperator{\supp}{supp}
\DeclareMathOperator{\argmax}{argmax}
\DeclareMathOperator{\argmin}{argmin}
\DeclareMathOperator{\E}{\mathbb{E}}
\csname endofdump\endcsname
\title{組合せ最適化特論}
\subtitle{第5回（最大流②）}
\author{担当: \textbf{相馬 輔}}
\date{2025/12/16}

\begin{document}
\maketitle

% \begin{frame}{}
    \centering\small
    \tikzset{block/.style={rectangle, draw=gray, fill=gray!10, thick, minimum width=3em, minimum height=2em, align=center}}
    \begin{tikzpicture}
        \node[block,] (BM) {重みなし\\二部マッチング};
        \node[block,above=1em of BM.north west, xshift=-2em] (WBM) {重みつき\\二部マッチング};
        \node[block,above=5em of BM.north east, xshift=1em] (NF) {最大流};
        \node[block,above=9em of BM] (MCF) {最小費用流};
        % \node[block,above=1em of NF.north east, xshift=1em] (MC) {最小カット};
        \node[block,right=11em of BM] (MTR) {マトロイド};
        \node[block,above=7em of MTR] (MI) {マトロイド交差};
        \node[block,above=3em of MCF, xshift=7em] (SFM) {劣モジュラ最小化};
        \graph[mygraph]{
            (BM) -> {(WBM), (NF)} -> (MCF);
            {(MTR), (BM)} -> (MI);
            % (NF) -> (MC);
            {(MCF), (MI)} -> (SFM);
        };
        \begin{scope}[on background layer]
            \node[mysubset, fit=(BM) (WBM), "二部マッチング" {below, UniBlue}]{};
            \node[mysubset, AlertOrange, fill=AlertOrange!10, fit=(NF) (MCF), "ネットワークフロー" {above left, AlertOrange}]{};
            \node[mysubset, UniGreen, fill=UniGreen!10, fit=(MTR) (MI), "マトロイド" {below, UniGreen}]{};
        \end{scope}
    \end{tikzpicture}
\end{frame}

\begin{frame}[allowframebreaks]
    \frametitle{各回の内容（予定）} 
    \small
    \begin{enumerate}
        \item (4/23) イントロ+多面体的組合せ論（線形計画法の復習，整数多面体，完全単模行列）
        \item[] \textcolor{gray}{(4/30) 休み}
        \item (5/7) 二部マッチング①（Konig-Egervaryの定理，増加道アルゴリズム，ハンガリー法）
        \item (5/14) 二部マッチング②（最短路問題の復習，逐次最短路法と主双対法，最適性基準からの見方）
        \item (5/21) 最大流① （定式化，最大流最小カット定理，応用例）
        \item (5/28) 最大流②（残余ネットワーク，Ford--Fulkerson法，Edmonds--Karpのアルゴリズム）+ 最小費用流①（定式化，輸送問題，最大流との関係）
        \item[] \textcolor{gray}{(6/4) 休み}
        % \item[] \sout{(12/11) 最小カット（永持茨木のアルゴリズム，Kargerのアルゴリズム）}
        \item (6/11) 最小費用流②（逐次最短路法，容量スケーリング法）
        \item[] \textcolor{gray}{(6/18) 休み}
        \item[] \textcolor{gray}{(6/25) 休み}
        \item (7/2) マトロイド（定義と公理系，貪欲法，マトロイド多面体） 
        \item (7/9) 独立マッチング・マトロイド交差 （定義，応用例，Edmondsの最大最小定理，増加道アルゴリズム）
        \item (7/16) 劣モジュラ関数①（諸例，劣モジュラ基多面体，Lov\'asz拡張，劣モジュラ最小化）
        \item[] \textcolor{gray}{(7/23) 休み}
        \item (7/30) 劣モジュラ関数②（劣モジュラ最大化，近似アルゴリズム，貪欲法）
        \item (どこか) 精選トピック①
        \item (どこか) 精選トピック②
        \item (どこか) 精選トピック③
    \end{enumerate}

    \bigskip
    \begin{itemize}
        \item レポート課題を6月と期末に出す予定
        \item 精選トピックの日時は今後調整 
    \end{itemize}
\end{frame}

\begin{frame}<1>[label=outline]{目次}
\begin{tikzpicture}
    \tikzstyle{block} = [rectangle, text width=\textwidth, rounded corners, minimum height=4em]
    \tikzstyle{line} = [draw, -latex]

    \node[block,alt={<1,3->{fill=cyan!10}{fill=red!10}}] (combopt) {
        \begin{columns}
            \begin{column}{0.5\linewidth}
        \textbf{1. Ford--Fulkerson法}
            \end{column}
            \begin{column}{0.4\linewidth}
                \small
                \structure{・} 残余ネットワーク\\
                \structure{・} 最大流最小カット定理の証明 \\
                \structure{・} 反復回数 \\
            \end{column}
        \end{columns}
    };
    \node[block,alt={<1-2,4->{fill=cyan!10}{fill=red!10}}, below=.5em of combopt] (LP) {
        \begin{columns}
            \begin{column}{0.5\linewidth}
        \textbf{2. Edmonds--Karp法}
            \end{column}
            \begin{column}{0.4\linewidth}
                \small
                \structure{・} Edmonds--Karp法 \\
                \structure{・} 距離ラベルを用いた解析 \\
                \structure{・} まとめ
            \end{column}
        \end{columns}
    };
\end{tikzpicture}%
\end{frame}

%%%%%%%%%%%%%%%%%%%%%%%%%%%%%%%%%%%%%%%%%%%%%
\section{Ford--Fulkerson法}
\againframe<2|handout:0>{outline}
%%%%%%%%%%%%%%%%%%%%%%%%%%%%%%%%%%%%%%%%%%%%%

\begin{frame}{最大流問題（復習）}
    $G = (V, A)$: 有向グラフ 

    \begin{definition}
        $\varphi: A \to \R$に対し，その\alert{境界} $\partial \varphi: V \to \R$を
        \[
            (\partial\varphi)(i) := \sum_{a \in \Out(i)} \varphi(a) - \sum_{a \in \In(i)} \varphi(a)  \qquad (i \in V)
        \]
        と定義する．ここで，$\In(i)$, $\Out(i)$はそれぞれ頂点$i$に入る枝，出る枝の集合．

        {\hfill\footnotesize ※$\delta^+(i)$, $\delta^-(i)$と表記する本もある．}
    \end{definition}
\end{frame}

\begin{frame}{最大流問題（復習）}
    $G = (V, A)$: 有向グラフ, $u: A \to \R_+$: 枝容量関数 \\
    $s, t \in V$: 始点, 終点 

    \begin{definition}[$s$--$t$流]
        $\varphi: A \to \R$が\alert{$s$--$t$\ruby{流}{フロー}} $\defiff$
        \begin{itemize}
            \item $0 \leq \varphi(a) \leq u(a)$ \qquad ($a \in A$) \hfill \textbf{（容量制約）} 
            \item $(\partial\varphi)(i) = 0$ \qquad ($i \in V \setminus \{s, t\}$) \hfill \textbf{（流量保存則）}
        \end{itemize}
    \end{definition}

    \pause
    \begin{block}{最大流問題 (Maximum Flow Problem)}
        $s$--$t$流$\varphi$で\alert{流量} \[
            \val(\varphi) := (\partial\varphi)(s) = -(\partial\varphi)(t) 
        \]
        が最大のものを求めよ．
    \end{block}
\end{frame}

\tikzgraphsset{declare={flowexample}{%
    s[at={(0,0)}, as={$s$}];
    a[at={(2, 1)}];
    b[at={(2,-1)}];
    c[at={(4, 1)}];
    d[at={(4,-1)}];
    t[at={(6, 0)}, as={$t$}];
}}

% \begin{frame}[label=flow-example]{例}
%     \textcolor{UniBlue}{$\varphi$}: 流量 / $u$: 容量

%     \begin{center}
%     \begin{tikzpicture}[scale=1.6]
%         \graph[no placement, mygraph]{flowexample};
%         \foreach \i/\j/\f/\u in {
%             s/a/5/5,
%             s/b/2/3,
%             a/b/2/2,
%             a/c/3/3,
%             d/a/0/2,
%             b/d/4/4,
%             c/t/4/5,
%             d/c/1/1,
%             d/t/3/3
%         } {
%             \path[myedges,->] (\i) edge node[above, midway, sloped] {\small \textcolor{UniBlue}{$\f$} / $\u$} (\j);
%         }
%     \end{tikzpicture}
%     \end{center}

%     流量: $7$ （実は最大流）
% \end{frame}

% \begin{frame}{最大流問題（復習）}
%     \footnotesize
%     \begin{primal}
%         \setlength{\abovedisplayskip}{0pt}
%         \begin{alignat*}{3}
%             \text{max}   \quad & \sum_{a \in \Out(s)} \varphi(a) && \quad \\
%             \text{s.t.}  \quad & \sum_{a \in \Out(i)} \varphi(a)  - \sum_{a \in \In(i)} \varphi(a) = 0 && \quad (i \in V \setminus \{s, t\}) \\
%                          \quad & 0 \leq \varphi(a) \leq u(a) && \quad (a \in A)
%         \end{alignat*}
%     \end{primal}
%     \begin{dual}
%         \setlength{\abovedisplayskip}{0pt}
%         \only<1>{
%         \begin{alignat*}{3}
%             \text{min}  \quad & \sum_{a \in A} u(a) y(a)   && \\
%             \text{s.t.} \quad & x(j) \phantom{{}-{}x(i)} + y(a) &&\;\geq 1  \quad && (a=sj \in \Out(s)) \\
%                         \quad & \phantom{x(i) } -x(i) + y(a)    &&\;\geq 0  \quad && (a=it \in \In(t)) \\
%                         \quad & x(j) - x(i) + y(a)              &&\;\geq 0  \quad && (a=ij \in A \setminus \Out(s) \cup \In(t)) \\
%                         \quad & y(a)                            &&\;\geq 0  \quad && (a \in A)
%         \end{alignat*}
%         }
%         \only<2>{
%         \begin{alignat*}{3}
%             \min_{x \in \R^V} & \quad \sum_{a=ij \in A} u(a) \max\{x(i) - x(j), 0\}   \\
%             \text{s.t.} &\quad  x(s) = 1  \\
%                         &\quad  x(t) = 0 
%         \end{alignat*}
%         }
%     \end{dual}
% \end{frame}

\begin{frame}{最大流最小カット定理}
    \small
    \begin{columns}[T]
    \begin{column}{.7\textwidth}
    \begin{definition}
        \begin{itemize}
            \item $X \subseteq V$が\alert{$s$--$t$カット} $\defiff s \in X, t \not\in X$
            \item $s$--$t$カット$X$の\alert{容量}: $u(\Out(X)) = \sum_{a \in \Out(X)} u(a)$
        \end{itemize}
    \end{definition}
    \end{column}
    \begin{column}{.3\textwidth}
    \centering

    \begin{tikzpicture}[scale=0.6]
        \graph[no placement, mygraph]{
            s[at={(0,0)}, as={$s$}];
            a[at={(2, 1)}];
            b[at={(2,-1)}];
            c[at={(4, 1)}];
            d[at={(4,-1)}];
            t[at={(6, 0)}, as={$t$}];
            s -> {a, b} -> {c, d} -> t;
            a -> {b, d};
            c -> d;
            {[edges={mypath}]
                a -> {c, d};
                b -> d;
            }
        };
        \node[AlertOrange, above] at ($(a)!.5!(c)$) {$\Out(X)$};
        \begin{scope}[on background layer]
        \node[mysubset, fit=(s) (a) (b),"{$X$}" {below, UniBlue}] {};
        \end{scope}
    \end{tikzpicture}
        % \tikz{\graph[mygraph, no placement]{
        %     subgraph C_n [n=4, clockwise];
        %     5[at={(3,0)}] -- 6[at={(2,1)}] -- 7[at={(2,-1)}] -- ;
        %     {[edges={mypath}]
        %         {(1),(2)} -- (6);
        %         {(2),(3)} -- (7);
        %     };
        % };
    \end{column}
    \end{columns}

    \pause
    \bigskip
    \structure{弱双対性} 任意の$s$--$t$流$\varphi$と$s$--$t$カット$X$に対し，次が成り立つ．
    \[
        \val(\varphi) \leq u(\Out(X))
    \]

    \pause
    \bigskip
    \begin{alertblock}{最大流最小カット定理 (Ford--Fulkerson, 1956; Danzig--Fulkerson, 1956)}
        \[
            \max_{\text{$\varphi$:~$s$--$t$流}} \val(\varphi) = \min_{\text{$X$:~$s$--$t$カット}} u(\Out(X)) 
        \]
    \end{alertblock}
\end{frame}

\begin{frame}{残余ネットワーク}
    \small 
    \begin{definition}[残余ネットワーク]
        \setlength{\abovedisplayskip}{2pt}
        \setlength{\belowdisplayskip}{2pt}
        \begin{itemize}
            \item 
    $s$--$t$流$\varphi$に対し，\alert{残余ネットワーク} $G_\varphi = (V, A_\varphi)$を次のように定める: 
    \begin{align*}
        A_\varphi &:= \Aphiup \cup \Aphilow \\
        \Aphiup &:= \{ a \in A : \varphi(a) < u(a) \} \quad\text{\color{gray}\footnotesize \dots\dots\dots$\varphi$を増やせる枝} \\
        \Aphilow &:= \{ a^{-1} : a \in A,\; \varphi(a) > 0  \} \quad \text{\color{gray}\footnotesize \dots\dots\dots$\varphi$を減らせる枝}
    \end{align*}
    ここで$a^{-1}$は$a$の逆向き枝．

    \item 残余ネットワークの枝容量$c_\varphi: A_\varphi \to \R_+$を
    次のように定める:
    \begin{align*}
        c_\varphi(a) := \begin{cases}
            u(a) - \varphi(a) & (a \in \Aphiup) \\
            \varphi(a^{-1})        & (a \in \Aphilow)
        \end{cases}
    \quad\text{\textcolor{gray}{\footnotesize \dots\dots\dots$\varphi$を増減できる余裕量}}
    \end{align*}
        \end{itemize}
    \end{definition}
\end{frame}

\begin{frame}{例}
    \small
   \newcommand{\flowvals}{%
            s/a/5/5,
            s/b/2/3,
            a/b/2/2,
            a/c/3/3,
            d/a/0/2,
            b/d/4/4,
            c/t/4/5,
            d/c/1/1,
            d/t/3/3
    }
   \begin{columns}[T]
   \begin{column}{.5\textwidth}
    \begin{tikzpicture}[xscale=1.1, yscale=1.2]
        \graph[no placement, mygraph]{flowexample};
        \foreach \i/\j/\f/\u in \flowvals {
            \path[myedges,->] (\i) edge node[above, midway, sloped] {\small \textcolor{UniBlue}{$\f$} / $\u$} (\j);
        }
    \end{tikzpicture}

   \end{column}
   \begin{column}{.5\textwidth}
    \vspace{1em}
    \begin{tikzpicture}[xscale=1.1, yscale=1.2]
        \graph[no placement, mygraph]{flowexample};
        \foreach \i/\j/\f/\u in \flowvals {
            % forward edges
            \ifthenelse{\f < \u}{
            \path[UniBlue, myedges,->, transform canvas={yshift=2pt}] (\i) edge node[above, midway, sloped] {\small \pgfmathparse{int(subtract(\u, \f))} \textcolor{UniBlue}{$\pgfmathresult$} } (\j);}{}
            % \path[UniBlue, myedges,->] ($(\i)!1pt!90:(\j)$) edge node[above, midway, sloped] {\small \pgfmathparse{int(subtract(\u, \f))} \textcolor{UniBlue}{$\pgfmathresult$} } ($(\j)!1pt!-90:(\i)$);}{}
            % backward edges
            \ifthenelse{\f > 0}{
            \path[myedges,->,AlertOrange, transform canvas={yshift=-1pt}] (\j) edge node[below, midway, sloped] {\small \textcolor{AlertOrange}{$\f$} } (\i);}{}
            % \path[myedges,->,AlertOrange] ($(\j)!1pt!90:(\i)$) edge node[below, midway, sloped] {\small \textcolor{AlertOrange}{$\f$} } ($(\i)!1pt!-90:(\j)$);}{}
        }
    \end{tikzpicture}
   \end{column}
   \end{columns} 

   \bigskip
   \begin{columns}[T]
   \begin{column}{.5\textwidth}
    \centering
    元のネットワーク $G$\\
    \textcolor{UniBlue}{$\varphi$}: 流量 / $u$: 容量
   \end{column}
   \begin{column}{.5\textwidth}
    \centering
    残余ネットワーク $G_\varphi$ \\
    $\Aphiup$, $\Aphilow$
   \end{column}
   \end{columns}
\end{frame}


\begin{frame}{増加道}

    \begin{definition}
        残余ネットワーク上の$s$--$t$パスを$\varphi$に対する\alert{増加道}と呼ぶ．
    \end{definition}

    \bigskip
    \structure{観察} 増加道があると，それに沿って流量を増やせる．

    \begin{center}
        \small
    残余ネットワーク $G_\varphi$ \\
    $\Aphiup$, $\Aphilow$ \\
        \tikz[baseline=(1.base)]{\graph[mygraph, edges={UniBlue}, grow right=4em]{
            1[as={$s$}] -> 2 -> 3 ->[AlertOrange] 4 -> 5 ->[AlertOrange] 6 ->[AlertOrange] 7 -> 8[as={$t$}] ;
        }} 
        \\[2em]
        元のグラフ$G$ \\
        \tikz{\graph[mygraph, grow right=4em]{
            1[as={$s$}] ->["{$+\eps$}" {above, UniBlue}] 2 ->["{$+\eps$}" {above, UniBlue}] 3 <-["{$-\eps$}" {above, AlertOrange}] 4 ->["{$+\eps$}" {above, UniBlue}] 5 <-["{$-\eps$}" {above, AlertOrange}] 6 <-["{$-\eps$}" {above, AlertOrange}] 7 ->["{$+\eps$}" {above, UniBlue}] 8[as={$t$}] ;
        }} 
    \end{center}
\end{frame}

\begin{frame}{増加道に沿った流の更新}
    増加道$P$に対し，次の$\varphi^+$を考える:
    \[
        \varphi^+(a) := \begin{cases}
            \varphi(a) \textcolor{UniBlue}{{}+\eps} & (a \in P \cap \Aphiup) \\
            \varphi(a) \textcolor{AlertOrange}{{}-\eps} & (a^{-1} \in P \cap \Aphilow) \\
            \varphi(a)        & (\text{otherwise})
        \end{cases}
    \]
    ここで$\eps = \min_{a \in P} c_\varphi(a)$である．このとき，$\varphi^+$は$s$--$t$流であり，その流量は$\val(\varphi^+) = \val(\varphi) + \eps$．
\end{frame}

\begin{frame}{増加道による最大流の特徴づけ}
    \begin{lemma}
        \begin{itemize}
            \item $s$--$t$流$\varphi$が最大流 $\iff$ $\varphi$に対する増加道が存在しない． 
            \item $\varphi$に対する増加道が存在しないとき，$G_\varphi$において$s$から到達可能な頂点集合$s \in X \subseteq V \setminus t$は\highlight{$\val(\varphi) = u(\Out(X))$}を満たす．
        \end{itemize}
    \end{lemma}
    \bigskip
    cf. 弱双対性 \highlight{$\val(\varphi) \leq u(\Out(X))$} \\
    最大二部マッチングの増加道による特徴づけ（第2回）の一般化
\end{frame}


\begin{frame}{補題の証明}
    \small
    \structure{増加道が存在しない $\implies$ $\varphi$は最大流:} 

    $X$を$G_\varphi$において$s$から到達可能な頂点集合とする．
    すると，$G_\varphi$において$X$から出る枝は存在しない．

    \bigskip
    ゆえに，以下が成り立つ:
    \begin{itemize}
        \item $G$において$X$から出る枝$a \in \Out(X)$に対し，$\varphi(a) = u(a)$
        \item $G$において$X$に入る枝$a \in \In(X)$に対し，$\varphi(a) = 0$
    \end{itemize}

    \bigskip
    よって，
    \begin{align*}
        \val(\varphi) &= \sum_{a \in \Out(X)} \varphi(a) - \sum_{a \in \In(X)} \varphi(a) \\
        &= \sum_{a \in \Out(X)} u(a) - 0 = u(\Out(X)).
    \end{align*}
    弱双対が等号で成り立つので，$\varphi$は最大流． \qed
\end{frame}

\begin{frame}{Ford--Fulkerson法}
    \small
    補題から直ちに次のアルゴリズムが得られる．

    \begin{block}{Ford--Fulkerson法}
    \begin{algorithmic}[1]
        \STATE $\varphi \equiv 0$.
        \WHILE {残余ネットワーク$G_\varphi$に増加道$P$が存在}
            \STATE $\eps \leftarrow \min_{a \in P} c_\varphi(a)$ \textcolor{gray}{\footnotesize \dotfill $P$に沿って増加できる最大の$\eps$}
            \STATE $\varphi(a) \gets 
            \begin{cases}
                \varphi(a) \textcolor{UniBlue}{{}+\eps} & a \in P \cap \Aphiup \\
                \varphi(a) \textcolor{AlertOrange}{{}-\eps} & a^{-1} \in P \cap \Aphilow \\
                \varphi(a)        & \text{otherwise}
            \end{cases}
            $
        \ENDWHILE
        \STATE $X :=$ $G_\varphi$において$s$から到達可能な頂点集合
        \RETURN $\varphi, X$
    \end{algorithmic}
    \end{block}

    \begin{theorem}[Ford--Fulkerson, 1957]
        Ford--Fulkerson法は停止した場合，最大流と最小カットを返す．
    \end{theorem}
\end{frame}

\begin{frame}{例}
   \small

   \only<1>{
   \newcommand{\flowvals}{%}
        s/a/0/5,
        s/b/0/3,
        a/b/0/2,
        a/c/0/3,
        d/a/0/2,
        b/d/0/4,
        c/t/0/5,
        d/c/0/1,
        d/t/0/3
   }}
    \only<2>{
   \newcommand{\flowvals}{%}
        s/a/0/5,
        s/b/2/3,
        a/b/0/2,
        a/c/2/3,
        d/a/2/2,
        b/d/2/4,
        c/t/2/5,
        d/c/0/1,
        d/t/0/3
   }}
    \only<3>{
   \newcommand{\flowvals}{%}
        s/a/2/5,
        s/b/2/3,
        a/b/2/2,
        a/c/2/3,
        d/a/2/2,
        b/d/4/4,
        c/t/2/5,
        d/c/0/1,
        d/t/2/3
   }}
    \only<4>{
   \newcommand{\flowvals}{%}
        s/a/3/5,
        s/b/2/3,
        a/b/2/2,
        a/c/3/3,
        d/a/2/2,
        b/d/4/4,
        c/t/3/5,
        d/c/0/1,
        d/t/2/3
   }}
    \only<5>{
   \newcommand{\flowvals}{%}
        s/a/4/5,
        s/b/2/3,
        a/b/2/2,
        a/c/3/3,
        d/a/1/2,
        b/d/4/4,
        c/t/3/5,
        d/c/0/1,
        d/t/3/3
   }}
    \only<6>{
   \newcommand{\flowvals}{%}
        s/a/5/5,
        s/b/2/3,
        a/b/2/2,
        a/c/3/3,
        d/a/0/2,
        b/d/4/4,
        c/t/4/5,
        d/c/1/1,
        d/t/3/3
   }}


   \begin{columns}[T]
   \begin{column}{.5\textwidth}
    \vspace{.9em}
    \begin{tikzpicture}[xscale=1.1, yscale=1.2]
        \graph[no placement, mygraph, math nodes]{flowexample};
        \foreach \i/\j/\f/\u in \flowvals
            \path[myedges,->] (\i) edge node[above, midway, sloped] {\small \textcolor{UniBlue}{$\f$} / $\u$} (\j);
    \end{tikzpicture}
   \end{column}
   \begin{column}{.5\textwidth}
    \vspace{1em}
    \begin{tikzpicture}[xscale=1.1, yscale=1.2]
        \node[inner sep=0pt, text width=1pt] at (0,1.45) {};
        \node[inner sep=0pt, text width=1pt] at (0,-1.45) {};
        \graph[no placement, mygraph, math nodes]{flowexample};
        \foreach \i/\j/\f/\u in \flowvals {
            % \def\opac{\ifthenelse{\f < \u}{1}{0}}
            % \path[UniBlue, myedges,->, transform canvas={yshift=2pt}, draw opacity=\opac] (\i) edge node[above, midway, sloped] {\small \pgfmathparse{int(subtract(\u, \f))} \textcolor{UniBlue}{$\pgfmathresult$} } (\j);
            % forward edges
            \ifthenelse{\f < \u}{
            \path[UniBlue, myedges,->, transform canvas={yshift=2pt}] (\i) edge node[above, midway, sloped] {\small \pgfmathparse{int(subtract(\u, \f))} \textcolor{UniBlue}{$\pgfmathresult$} } (\j);}{}
            % backward edges
            \ifthenelse{\f > 0}{
            \path[myedges,->,AlertOrange, transform canvas={yshift=-1pt}] (\j) edge node[below, midway, sloped] {\small \textcolor{AlertOrange}{$\f$} } (\i);}{}
            % \def\opac{\ifthenelse{\f > 0}{1}{0}}
            % \path[myedges,->,AlertOrange, transform canvas={yshift=-1pt}, draw opacity=\opac] (\j) edge node[below, midway, sloped] {\small \textcolor{AlertOrange}{$\f$} } (\i);
        }

        \begin{scope}[on background layer]
            \draw<1>[line width=1em, yellow] (s) -- (b) -- (d) -- (a) -- (c) -- (t);
            \draw<2>[line width=1em, yellow] (s) -- (a) -- (b) -- (d) -- (t);
            \draw<3>[line width=1em, yellow] (s) -- (a) -- (c) -- (t);
            \draw<4>[line width=1em, yellow] (s) -- (a) -- (d) -- (t);
            \draw<5>[line width=1em, yellow] (s) -- (a) -- (d) -- (c) -- (t);
        \end{scope}
    \end{tikzpicture}
   \end{column}
   \end{columns} 

   \bigskip
   \begin{columns}[T]
   \begin{column}{.5\textwidth}
    \centering
    元のネットワーク $G$\\
    \textcolor{UniBlue}{$\varphi$}: 流量 / $u$: 容量
   \end{column}
   \begin{column}{.5\textwidth}
    \centering
    残余ネットワーク $G_\varphi$ \\
    $\Aphiup$, $\Aphilow$
   \end{column}
   \end{columns}
   \bigskip
   流量: \textcolor{UniBlue}{\only<1>{$0$}\only<2>{$2$}\only<3>{$4$}\only<4>{$5$}\only<5>{$6$}\only<6>{$7$}}
   \hfill \onslide<6>{【終了】}
\end{frame}

\begin{frame}{Ford--Fulkerson法}
     残余ネットワークの構築・増加道の探索・流の更新: $O(m)$時間 \\
     {\footnotesize \hfill ($m = |A|$)}

    \bigskip
    \structure{反復回数}
    \begin{itemize}
        \setlength\itemsep{1em}

        \item \alert{枝容量が整数値}のとき: 各反復で流量が少なくとも$1$増加．したがって，最大流量を$F$とすると，高々$F$回の反復．\\ 
        $\longrightarrow$ 全体では$O(m F)$時間 （\alert{擬多項式時間}）

        \item 枝容量が無理数値のとき: 増加道の選び方によっては，有限回で停止せず，かつ最大でない流に収束してしまう例がある \mycite{Ford--Fulkerson,~1962}
    \end{itemize}
\end{frame}

%%%%%%%%%%%%%%%%%%%%%%%%%%%%%%%%%%%%%%%%%%%%%
\section{Edmonds--Karp法}
\againframe<3|handout:0>{outline}
%%%%%%%%%%%%%%%%%%%%%%%%%%%%%%%%%%%%%%%%%%%%%
\begin{frame}{Edmonds--Karp法}
    Ford--Fulkerson法において，\alert{枝数が最小の増加道}を選ぶ． 
    
    \bigskip
    たとえば，幅優先探索で$s$--$t$パスを探索すればよい．

    \bigskip
    \begin{theorem}[Edmonds--Karp, 1972]
        \textbf{任意の}枝容量に対し，Edmonds--Karp法は$O(mn)$回の反復で終了する．したがって，$O(m^2 n)$時間で最大流と最小カットを求められる．
    \end{theorem}
    \hfill \alert{強多項式時間} \\
    {\footnotesize \hfill ($n = |V|$, $m = |A|$)} 
\end{frame}

\begin{frame}{Edmonds--Karp法の解析}
    $\varphi$: Edmonds--Karp法の途中の流 \\
    $P$: 選ばれた$\varphi$増加道 \\
    $\varphi^+$: $P$に沿って更新した流

    \vfill
    \begin{lemma}[距離ラベルの単調性]
        $d(i) :=$ $G_\varphi$における頂点$i$から$t$までの最短距離（枝数） \\
        $d^+(i) :=$ $G_{\varphi^+}$における頂点$i$から$t$までの最短距離（枝数） \\
        とする．\\ 
        このとき，任意の頂点$i \in V$に対し，$d(i) \leq d^+(i)$が成り立つ．
    \end{lemma}
\end{frame}

\begin{frame}{Edmonds--Karp法の解析}
    \structure{(証明)}
    背理法．$d(i) > d^+(i)$となる頂点$i$が存在すると仮定する．
    そのような頂点$i$の中で$d^+(i)$が最小のものを取る．{\footnotesize ($i \neq t$)}
    
    \bigskip
    $j$: $G_{\varphi^+}$において$i$から$t$までの最短路における$i$の直後の頂点 
    
    \bigskip
    \structure{Claim.} $ij \notin A_{\varphi}$ \\
    {\small \structure{(∵)} もし$ij \in A_{\varphi}$ならば，
    \[
        d(i) \leq d(j) + 1 \leq d^+(j) + 1 = d^+(i) < d(i)
    \]
    となり矛盾．\qed
    }

\end{frame}

\begin{frame}{Edmonds--Karp法の解析}
    \structure{Claim.} $ij \notin A_{\varphi}$ \\

    \bigskip
    したがって，
    \begin{itemize}
        \item $ji$が$G_\varphi$に存在し，その逆向き枝$ij$は存在しない．
        \item また，更新後の$G_{\varphi^+}$に$ij$が現れているので，$ji$は$G_\varphi$の最短路$P$に含まれていたことになる．よって，\highlight{$d(j) = d(i) + 1$}.
    \end{itemize}

    \bigskip
    すると，
    \begin{align*}
        d(i) &= d(j) - 1  \tag{上式を移項}\\ 
             &\leq d^+(j) - 1 \tag{$d(j) \leq d^+(j)$} \\
             &= d^+(i) - 2 \tag{$j$は$G_{\varphi^+}$において$i$の直後の頂点}
    \end{align*}
    となり，そもそもの仮定$d(i) > d^+(i)$に矛盾．
    \qed
\end{frame}

% \begin{frame}{Edmonds--Karp法の解析}
%     \begin{lemma}
%         ある更新で枝$ij$が最短路に含まれたたとする．
%         さらに，その後の更新で$ji$が最短路に含まれたとする．
%         それぞれの時点における距離ラベルを$d, d'$とすると，$d'(i) \geq d(i) + 2$が成り立つ．
%     \end{lemma}
%     \structure{(証明)} $uv$が最短路に含まれるとき，$d(u) = d(v) + 1$が成り立つことを用いる．

%     TODO: 図を描く
%     \qed
% \end{frame}

\begin{frame}{Edmonds--Karp法の解析}
    \small

    \begin{theorem}[Edmonds--Karp]
        任意の枝容量に対し，Edmonds--Karp法は$O(mn)$回の反復で終了する．したがって，$O(m^2 n)$時間で最大流と最小カットを求められる．
    \end{theorem}


    \structure{(証明)}
    $ij \in A_\varphi$が，更新後に$ij \notin A_{\varphi^+}$となったとき，$ij$は\alert{飽和}したという．

    \begin{itemize}[<+->]
        \setlength\itemsep{0em}
        \item 各反復では，$P$の少なくとも1本の枝が飽和する．飽和した枝$ij$について$d(i) = d(j) + 1$が成り立つ．
        % $\longrightarrow$ 高々$m$回の反復後に，$ij$と$ji$の両方が最短路に含まれる． \\
        % $\longrightarrow$ 高々$m$回の反復後に，ある$i$で$d(i)$が$2$以上増大する（補題）

        \item 一度飽和した$ij$が再び飽和するには，その間に$ji$が一度は最短路に含まれる必要がある．そのときの距離ラベルに関して，$d(j) = d(i) + 1$が成り立つ．
        
        \item 距離ラベル$d(i) \leq n$ or $d(i) = +\infty$ ($i$から$t$へ到達不能の場合) 
        % $\longrightarrow$ 各$d(i)$は高々$n$回増大しうる
    \end{itemize}
    \bigskip
    \onslide<+->
    よって，距離ラベルの単調性と合わせると，各枝は高々$n/2$回飽和しうる．ゆえに全体の反復は$O(mn)$回． 
    \qed
\end{frame}

\begin{frame}{Edmonds--Karp法の解析}
    \begin{center}
        \begin{tikzpicture}
            %grid
            \draw[gray!30] (0,0) grid (11, 5);
            \draw[->, thick] (0,0) -- (0, 5);
            \node[left] at (0,0) {$0$};
            \node[left] at (0,5) {$n$};

            \tikzset{mydot/.style={circle, inner sep=1.5pt}}
            \node[mydot, fill=UniBlue, "{$d(i)$}" below]  (di)   at (1 ,2) {};
            \node[mydot, fill=UniPink, "{$d(j)$}" below]  (dj)   at (2 ,1) {};
            \node[mydot, fill=UniBlue, "{$d'(i)$}" below] (dpi)  at (5 ,2) {};
            \node[mydot, fill=UniPink, "{$d'(j)$}" below] (dpj)  at (6 ,3) {};
            \node[mydot, fill=UniBlue, "{$d''(i)$}" below] (dppi) at (9 ,4) {};
            \node[mydot, fill=UniPink, "{$d''(j)$}" below] (dppj) at (10,3) {};
            \node[align=center, font=\footnotesize] at (di |- 0,-1) {$ij$が飽和したとき \\ $d(i) = d(j) + 1$};
            \node[align=center, font=\footnotesize] at (dpj |- 0,-1) {$ji$が最短路に含まれたとき \\ $d'(j) = d'(i) + 1$};
            \node[align=center, font=\footnotesize] at (dppj |- 0,-1) {$ij$が再度飽和したとき \\ $d''(i) = d''(j) + 1$};
        \end{tikzpicture}
    \end{center}
\end{frame}

\begin{frame}{まとめ}
    \structure{最大流問題}
    \begin{itemize}
        \item 最大流最小カット定理，整数性
        \item Ford--Fulkerson法: 枝容量が整数値のとき$O(m F)$時間
        \item Edmonds--Karp法: $O(m^2 n)$時間
    \end{itemize}

    \bigskip
    \structure{その他の最大流アルゴリズム}
    \begin{itemize}
        \setlength{\itemsep}{1em}
        \item Dinitz (1970): $O(mn^2)$時間 \\
        動的木を用いると $O(mn \log n)$時間 (Sleator--Tarjan, 1983)
        \item プッシュ・再ラベル (Goldberg--Tarjan, 1988): $O(n^2\sqrt{m})$時間 \\
        動的木を用いると $O(mn \log (n^2/m))$時間 
    \end{itemize}
    
\end{frame}

\end{document}
